\documentclass[11pt,letterpaper]{article}
\usepackage[T1]{fontenc}
\usepackage[utf8]{inputenc}
\usepackage{newtxtext}
\usepackage{amsmath,amssymb,amsthm,mathtools}
\usepackage{newtxmath}
\usepackage[margin=1.03in,headheight=15pt]{geometry}
\usepackage{microtype}
\usepackage{xcolor,booktabs,array,enumitem}
\usepackage[most]{tcolorbox}
\usepackage{titlesec,fancyhdr}
\usepackage{hyperref}
\usepackage{bookmark}
\definecolor{accent}{HTML}{214F47}
\definecolor{soft}{HTML}{F0F5F2}
\definecolor{muted}{HTML}{536561}
\hypersetup{colorlinks=true,linkcolor=accent,citecolor=accent,urlcolor=accent,
  pdftitle={A two-valued obstruction to gluing Frechet powers},
  pdfauthor={Research manuscript prepared with ChatGPT},
  pdfsubject={A proposed solution of Massas's sheaf question; complete proofs}}
\titleformat{\section}{\Large\bfseries\color{accent}}{\thesection}{0.7em}{}
\titleformat{\subsection}{\large\bfseries\color{accent}}{\thesubsection}{0.7em}{}
\pagestyle{fancy}
\fancyhf{}
\fancyhead[L]{\small\color{muted}A two-valued obstruction to gluing Fr\'echet powers}
\fancyhead[R]{\small\color{muted}Research manuscript}
\fancyfoot[C]{\small\thepage}
\renewcommand{\headrulewidth}{0.3pt}
\setlength{\parskip}{0.3em}
\setlength{\parindent}{1.2em}
\setlength{\emergencystretch}{2em}
\setlist{itemsep=0.2em,topsep=0.4em}
\newtheorem{theorem}{Theorem}[section]
\newtheorem{lemma}[theorem]{Lemma}
\newtheorem{proposition}[theorem]{Proposition}
\newtheorem{corollary}[theorem]{Corollary}
\theoremstyle{definition}
\newtheorem{definition}[theorem]{Definition}
\newtheorem{example}[theorem]{Example}
\theoremstyle{remark}
\newtheorem{remark}[theorem]{Remark}
\newcommand{\om}{\omega}
\newcommand{\Fr}{\mathrm{Fr}}
\newcommand{\cF}{\mathcal F}
\newcommand{\RA}{\mathscr R_A}
\newcommand{\SA}{\mathscr S_A}
\newcommand{\Pow}{\mathcal P}
\newcommand{\Ult}{\operatorname{Ult}}
\newcommand{\gen}[1]{\langle #1\rangle}
\newcommand{\eqset}[2]{E(#1,#2)}
\newcommand{\co}{\mathrm c}
\newcommand{\const}[1]{\overline{#1}}
\newcommand{\ZF}{\mathsf{ZF}}
\newcommand{\DC}{\mathsf{DC}}
\newcommand{\ZFC}{\mathsf{ZFC}}
\newcommand{\B}{\mathbb B}
\newcommand{\I}{\mathcal I}
\newcommand{\res}[2]{#1\!\upharpoonright_{#2}}
\newtcolorbox{resultbox}{colback=soft,colframe=accent,boxrule=0.6pt,arc=1mm,
  left=3mm,right=3mm,top=2mm,bottom=2mm,breakable}

\begin{document}
\begin{titlepage}
\thispagestyle{empty}
\vspace*{0.6in}
{\sffamily\color{accent}\large NONSTANDARD ANALYSIS \; / \; FILTERS \; / \; SHEAF THEORY\par}
\vspace{0.45in}
{\Huge\bfseries A two-valued obstruction\par to gluing Fr\'echet powers\par}
\vspace{0.3in}
{\LARGE\color{muted}A proposed solution of Massas's sheaf question\par}
\vspace{0.5in}
{\large Complete proofs, local classification, and reproducible checks\par}
\vspace{0.25in}
{\large 20 September 2026\par}
\vspace{0.45in}
\begin{resultbox}
\textbf{Main result, in $\ZF$.} On the category of all proper filters on
$\om$ extending the cofinite filter, equipped with the dense topology,
\[
  F\longmapsto A^\om/{\sim_F}
  \quad\text{is a sheaf if and only if}\quad |A|\leq 1.
\]
Two distinct values already suffice for failure. An explicit countably
generated covering sieve carries compatible constant local sections with
no global representative.
\end{resultbox}
\vfill
\noindent\textbf{Research status.} This is an unrefereed manuscript prepared
with ChatGPT. It presents a complete mathematical argument, not a claim of
independent expert certification. The motivating question is explicitly open
in the cited 2023 article and 2024 dissertation. A targeted literature check
found no later resolution, but priority and novelty have not been established.
The computational supplements check exact finite consequences; they are not a
formalization of the infinite proof.
\end{titlepage}

\begin{abstract}
Massas asks whether uncountability is necessary for the reduced-power
presheaf associated with a structure to fail the sheaf condition on the
Fr\'echet category. We give a negative answer using only two distinct elements.
Partition $\om$ into infinitely many infinite rows. The filters concentrated
on individual rows, together with the filter concentrated on tails of rows,
generate a dense covering sieve. Assign one constant value on every row
component and a different constant value on the tail component. These
sections match because the components are incompatible. A global sequence
would nevertheless have to take both values cofinitely on one infinite row.
The construction and the resulting classification are proved in $\ZF$.
We also prove separatedness, a finite-pasting theorem showing that the
countable number of generators is sharp, and failure above every countably
generated free filter without any choice axiom. Under dependent choice,
sheafness above an arbitrary base filter is equivalent, for a nontrivial
value set, to finiteness of its quotient Boolean algebra. Finally, in
$\ZFC$, we identify the associated sheaf with a product of ultrapowers and
explain the obstruction as a missing two-valued section. Every extension is
proved independently of the literature, with its foundational assumptions
stated explicitly.
\end{abstract}

\tableofcontents
\newpage

\section{The question and the result}
\label{sec:question}

\subsection{A specific published open problem}
In his treatment of hyperreal numbers by possibility semantics, Massas
considers a varying reduced power indexed by proper filters extending the
cofinite filter on $\om$. In \cite[Definitions 5.4--5.7 and Fact 5.10]{Massas2023},
he proves that this presheaf need not be a sheaf. His counterexample assumes
an injection $2^\om\hookrightarrow A$; with choice, he explains that the
argument works for every uncountable domain. The question immediately after
Fact~5.10 asks whether uncountability is necessary. It appears on printed
page~517. The same question is repeated after Fact~7.4.10, on printed
page~249 of his 2024 dissertation \cite{Massas2024}.

The issue concerns the \emph{underlying set presheaf}. We write it as
$\RA$ to distinguish it from a single ultrapower or a single ordered field.
Its sections at a filter $F$ are sequences in $A$, identified when they
agree on an $F$-large set:
\begin{equation}\label{eq:presheaf-intro}
 \RA(F)=A^\om/{\sim_F},\qquad
 f\sim_F g\ \Longleftrightarrow\ \{j:f(j)=g(j)\}\in F.
\end{equation}
A sheaf condition asks whether compatible sections given at enough stronger
filters come from one section at the original filter.

\begin{theorem}[Global classification in $\ZF$]\label{thm:main}
For every set $A$, the presheaf $\RA$ on the Fr\'echet category with its
dense topology is a sheaf if and only if $A$ has at most one element.
If $a_0,a_1\in A$ are distinct, failure occurs at the cofinite filter on a
covering sieve generated by countably many explicitly defined, countably
generated filters. All local sections in the counterexample are represented
by one of the two constant sequences $\const{a_0}$ and $\const{a_1}$.
\end{theorem}

The theorem answers the selected question negatively and supplies the exact
cardinality boundary. For the usual convention that a first-order structure
has a nonempty domain, the boundary is simply that the domain is a
singleton. The proof is given in Section~\ref{sec:obstruction}; no assertion
about unpublished work or priority is needed for its validity.

\subsection{Why the construction works}
The key distinction is between being compatible as filters and imposing
requirements on overlapping subsets of $\om$. A filter on one row and a
filter on tails of rows have no common proper extension. Consequently, a
matching family is allowed to assign them different constant values.
Nevertheless, any \emph{single sequence} representing both sections would
have to be constant with two different values on almost all of a suitable
row. This is impossible.

The argument replaces the large family of different values used in the
published cardinality obstruction by two values and a tail condition. The
obstruction is therefore not a shortage of enough distinct values in the
domain. It is a failure of countable patching for representatives of local
equivalence classes.

\subsection{Scope and foundations}
Sections~\ref{sec:site}--\ref{sec:countablebase} work in classical $\ZF$,
with no axiom of choice. ``Choice-free'' here does not mean intuitionistic:
ordinary classical case distinctions are used. Section~\ref{sec:boolean}
uses dependent choice only for the infinite direction of a local Boolean
classification. Section~\ref{sec:sheafification} is explicitly in $\ZFC$
and uses ultrafilter extensions. Neither of those later assumptions enters
the main theorem.

The article does not propose a new arithmetic of all hyperreal or surreal
numbers, and does not address an ordinal-arithmetic conjecture. It solves a
particular gluing question that arises in a construction of the hyperreal
line. The mathematical claims below are proved in full; novelty remains a
separate bibliographical issue.

\section{The filter site and its reduced powers}
\label{sec:site}

\subsection{Filters and positive sets}
Write $\om=\{0,1,2,\ldots\}$ and $S^\co=\om\setminus S$.
A \emph{proper filter} on $\om$ is a family $F\subseteq\Pow(\om)$ that
contains $\om$, omits $\varnothing$, is closed under finite intersections,
and is upward closed under inclusion. The cofinite filter is
\[
 \Fr=\{X\subseteq\om:\om\setminus X\text{ is finite}\}.
\]
Throughout, a \emph{free filter} means a proper filter containing $\Fr$.
Such a filter has no finite member: a finite member and its cofinite
complement would have empty intersection.

Given a family of subsets with the finite-intersection property,
$\gen{\mathcal E}$ denotes the filter it generates. The same notation is
used for an improper generated filter when discussing compatibility, but
an improper filter is never an object of our site.

\begin{definition}
For a proper filter $F$, a set $S\subseteq\om$ is \emph{$F$-positive} if
$S^\co\notin F$. Two free filters $F,G$ are \emph{compatible} if some
proper free filter contains both; otherwise they are incompatible.
\end{definition}

\begin{lemma}[Adjoining a set]\label{lem:positive}
Let $F$ be a proper filter and $S\subseteq\om$. Then
\[
 \gen{F\cup\{S\}}\text{ is proper}\quad\Longleftrightarrow\quad S^\co\notin F.
\]
When this generated filter is proper, its members are precisely the sets
$X$ for which $C\cap S\subseteq X$ for some $C\in F$. If $F$ is free,
$S$ is $F$-positive if and only if $C\cap S$ is infinite for every $C\in F$.
\end{lemma}
\begin{proof}
Finite intersections of generators have the form $C$ or $C\cap S$ with
$C\in F$. The latter is empty for some $C$ exactly when some member of
$F$ is contained in $S^\co$, equivalently when $S^\co\in F$. This proves
the first assertion and the description of the generated filter.

Now suppose $F$ is free and $C\cap S$ is finite. Intersecting $C$ with
the cofinite complement of $C\cap S$ produces an $F$-member disjoint
from $S$, so $S^\co\in F$. Conversely, if $S^\co\in F$, its intersection
with $S$ is empty. These implications prove the final assertion.
\end{proof}

Two filters are incompatible exactly when there exist $C\in F$ and
$D\in G$ with $C\cap D=\varnothing$. Indeed, such an intersection
prevents a proper common extension; if no such intersection is empty,
$\gen{F\cup G}$ itself is a proper common extension.

\subsection{Objects, arrows, and covering sieves}
Let $\cF$ be the set of free filters on $\om$. Regard it as a category
with a unique arrow $G\to F$ when $G\supseteq F$. Thus a stronger filter
is the source of an arrow to a weaker one. For clarity we use actual set
inclusion rather than a separate reversed-order symbol.

A \emph{sieve at $F$} is a family $D\subseteq\cF$ such that every
$G\in D$ extends $F$, and every extension of a member of $D$ is also
in $D$. It is \emph{covering} when
\begin{equation}\label{eq:density}
 \text{for every }G\supseteq F\text{ in }\cF,
 \text{ there is }H\supseteq G\text{ with }H\in D.
\end{equation}
These are exactly the covers used in the question. The sieve formulation
is also the standard language for Grothendieck topologies
\cite[Section 7.48]{StacksTopology}.

For $H\in\cF$, write
\[
 \uparrow H=\{G\in\cF:H\subseteq G\}.
\]
Here the arrow points upward in \emph{inclusion}, not upward in the
opposite categorical order. A sieve \emph{generated by} filters $H_i$
is $\bigcup_i\uparrow H_i$. A countably generated sieve need not be a
countable set of filters. Separately, a filter is \emph{countably generated}
when it has a countable filter base. We will use both notions and keep
them distinct.

\begin{lemma}\label{lem:topology}
The covering sieves just defined form a Grothendieck topology.
Every covering sieve is nonempty.
\end{lemma}
\begin{proof}
The maximal sieve of all extensions of $F$ is dense. If $D$ is a covering
sieve at $F$ and $G\supseteq F$, the sieve of members of $D$ extending
$G$ is dense above $G$: apply \eqref{eq:density} to an arbitrary
extension of $G$.

For transitivity, suppose $D$ covers $F$ and $E$ is a sieve at $F$ such
that $E\cap\uparrow G$ covers $G$ for each $G\in D$. Given $H\supseteq F$,
choose an extension $G\supseteq H$ in $D$, then an extension
$K\supseteq G$ in $E$. Thus $E$ is dense. These are the maximality,
pullback, and transitivity axioms. Applying density to $F$ itself shows
that a covering sieve cannot be empty. All choices in this argument are
within a fixed finite chain of existential assertions.
\end{proof}

\subsection{Sections and matching families}
For sequences $f,g:\om\to A$, put
\[
 \eqset{f}{g}=\{j\in\om:f(j)=g(j)\}.
\]
The relation in \eqref{eq:presheaf-intro} is an equivalence relation.
If $G\supseteq F$, equality modulo $F$ implies equality modulo $G$;
hence there is a well-defined restriction map
\[
 \rho_{F,G}:\RA(F)\longrightarrow\RA(G),\qquad
 \rho_{F,G}([f]_F)=[f]_G.
\]
Identity and composition hold directly, so $\RA$ is a presheaf.

A \emph{matching family} on a sieve $D$ is a collection
$s_G\in\RA(G)$, one for each $G\in D$, such that
\[
 \rho_{G,H}(s_G)=s_H\quad\text{whenever }G,H\in D\text{ and }H\supseteq G.
\]
An \emph{amalgam} at $F$ is $s\in\RA(F)$ with
$\rho_{F,G}(s)=s_G$ for every $G\in D$. The presheaf is a \emph{sheaf}
if every matching family on every covering sieve has exactly one amalgam.
It is \emph{separated} if it has at most one.

There is an equivalent way to prescribe data on generators. Given filters
$H_i$, sections at them determine a matching family on the generated sieve
provided they agree at every common proper extension. If a filter extends
several generators, this agreement makes the definition independent of
which generator is used. No agreement condition is imposed between
incompatible generators.

\section{The explicit two-valued obstruction}
\label{sec:obstruction}

\subsection{Rows and their tail filter}
Fix an explicit partition
\begin{equation}\label{eq:partition}
 \om=\bigsqcup_{n<\om}B_n
\end{equation}
into infinite sets, and let
\begin{equation}\label{eq:tail}
 T_m=\bigcup_{n\geq m}B_n
     =\om\setminus\bigcup_{n<m}B_n.
\end{equation}
An arithmetic formula for the partition is supplied in
Section~\ref{sec:arithmetic}; therefore no choice is hidden in its existence.
Define the row filters and the tail filter by
\begin{align}
 F_n&=\{X\subseteq\om:B_n\setminus X\text{ is finite}\},\label{eq:rowfilter}\\
 F_\infty&=\{X\subseteq\om:\text{for some }m,
                         \ T_m\setminus X\text{ is finite}\}.\label{eq:tailfilter}
\end{align}
Equivalently,
\[
 F_n=\gen{\Fr\cup\{B_n\}},\qquad
 F_\infty=\gen{\Fr\cup\{T_m:m<\om\}}.
\]

\begin{lemma}\label{lem:incompatible}
All these filters are proper and free. They are countably generated.
The filters $F_n$ are pairwise incompatible, and every $F_n$ is
incompatible with $F_\infty$.
\end{lemma}
\begin{proof}
The defining families are upward closed and contain every cofinite set.
Closure under intersections for $F_\infty$ follows by taking the larger
of two tail indices and the union of the two finite exceptional sets.
No row or tail is finite, so neither definition includes the empty set.
A countable base for $F_n$ is
$\{B_n\setminus\{0,\ldots,r-1\}:r<\om\}$; a countable base for
$F_\infty$ uses $T_m\setminus\{0,\ldots,r-1\}$ for $(m,r)\in\om^2$.

For $n\neq k$, the disjoint sets $B_n\in F_n$ and $B_k\in F_k$
preclude a proper common extension. Also $B_n\in F_n$ and
$T_{n+1}\in F_\infty$ are disjoint. This proves every claimed
incompatibility.
\end{proof}

\subsection{A countably generated dense sieve}
Consider
\begin{equation}\label{eq:D}
 D=\uparrow F_\infty\ \cup\ \bigcup_{n<\om}\uparrow F_n.
\end{equation}
It is a sieve at $\Fr$, and its displayed components are disjoint.

\begin{lemma}[Density]\label{lem:dense}
The sieve $D$ covers $\Fr$.
\end{lemma}
\begin{proof}
Let $G\supseteq\Fr$ be any proper free filter. If some row $B_n$ is
$G$-positive, Lemma~\ref{lem:positive} says that
\[
 H=\gen{G\cup\{B_n\}}
\]
is a proper free filter. It extends both $G$ and $F_n$, so $H\in D$.

Otherwise $B_n^\co\in G$ for every $n$. Since $G$ is closed under
\emph{finite} intersections, for each $m$ it contains
\[
 \bigcap_{n<m}B_n^\co=T_m.
\]
It consequently contains $F_\infty$, and $G$ itself belongs to $D$.
Both cases provide the extension required by \eqref{eq:density}.
\end{proof}

The second case does not use countable completeness of $G$. For a fixed
$m$, only $m$ sets are intersected. Nor does the first case require a
simultaneous choice of positive rows for all filters: the proof establishes
an extension for one arbitrary $G$ at a time.

\subsection{A matching family with no amalgam}
Fix distinct $a_0,a_1\in A$, and let $\const{a_i}$ be the constant
sequence with value $a_i$. For each $G\in D$, define
\begin{equation}\label{eq:matching}
 s_G=
 \begin{cases}
  [\const{a_0}]_G,&\text{if }G\supseteq F_n\text{ for some }n,\\
  [\const{a_1}]_G,&\text{if }G\supseteq F_\infty.
 \end{cases}
\end{equation}
Lemma~\ref{lem:incompatible} makes the cases disjoint. Every extension
of a filter in a component remains in that component, so restriction
preserves its assigned constant section. Thus \eqref{eq:matching} is a
matching family.

\begin{proposition}[The contradiction]\label{prop:contradiction}
The matching family \eqref{eq:matching} has no amalgam in $\RA(\Fr)$.
\end{proposition}
\begin{proof}
Suppose $[f]_{\Fr}$ were an amalgam. Its restriction to the tail filter
would satisfy
\[
 E_1:=\{j:f(j)=a_1\}\in F_\infty.
\]
By \eqref{eq:tailfilter}, there is $m$ such that $T_m\setminus E_1$ is
finite. Its restriction to the row filter $F_m$ would satisfy
\[
 E_0:=\{j:f(j)=a_0\}\in F_m,
\]
so $B_m\setminus E_0$ is finite. Because $B_m\subseteq T_m$, the set
$B_m\setminus E_1$ is also finite. The infinite set $B_m$ therefore
has an element in $E_0\cap E_1$. At that element, $f$ has both values
$a_0$ and $a_1$, contrary to their being distinct.
\end{proof}

\begin{proof}[Proof of Theorem~\ref{thm:main}]
If $A$ has two distinct elements, the preceding construction proves
failure of sheafness. If $A$ is a singleton, $\RA(F)$ is a singleton
for every object and every matching family has one amalgam. If
$A=\varnothing$, every $\RA(F)$ is empty because $\om$ is nonempty.
Every covering sieve is nonempty, so it has no matching family with
values in this presheaf; the sheaf condition is then vacuously true.
These exhaust the cases in classical $\ZF$.
\end{proof}

\subsection{Why several tempting objections do not apply}
The family is not assigning different values on two compatible pieces.
The row and tail filters are incompatible because the tail filter contains
\emph{all} row tails, including the one disjoint from the chosen row. A
putative representative at the tail filter, however, agrees with its
constant value outside finitely many points of \emph{one} tail. Selecting
the index of that tail identifies a row on which the global contradiction
occurs.

The proof does not exchange two universal quantifiers, assume a uniform
exceptional set for all rows, or require that $f$ be two-valued everywhere.
It uses the tail index supplied by $f$ and the exceptional set supplied
for that single row. Other values of $f$, when $A$ is larger, cannot
remove the contradiction.

Finally, $D$ is the full extension-closed family in \eqref{eq:D}, not
just the listed generators. This distinction is essential for matching
the covering definition in the original question.

\section{An arithmetic realization and an exact witness}
\label{sec:arithmetic}

\subsection{A partition determined by powers of two}
Take
\begin{equation}\label{eq:dyadic-rows}
 B_n=\{2^n(2k+1)-1:k<\om\}.
\end{equation}
Every positive integer has a unique factorization as a power of two times
an odd integer. Applied to $j+1$, this shows that \eqref{eq:dyadic-rows}
is a partition of $\om$ into infinite sets. In this realization,
\begin{equation}\label{eq:dyadic-tails}
 T_m=\{j<\om:2^m\text{ divides }j+1\}.
\end{equation}
The first few sets are
\[
\begin{array}{ll}
 B_0=\{0,2,4,6,\ldots\}, & B_1=\{1,5,9,13,\ldots\},\\
 B_2=\{3,11,19,27,\ldots\}, & B_3=\{7,23,39,55,\ldots\}.
\end{array}
\]
One may view the rows as the finite $2$-adic valuations of $j+1$; no
$2$-adic analysis is needed.

\subsection{Turning cofinite bounds into a contradictory index}
Suppose the hypothetical amalgam provides integers $m,N_0,N_1$ such that
\begin{align*}
  j\in B_m,\ j\geq N_0&\quad\Longrightarrow\quad f(j)=a_0,\\
  j\in T_m,\ j\geq N_1&\quad\Longrightarrow\quad f(j)=a_1.
\end{align*}
Let $M=\max(N_0,N_1)$, $r=2^m-1$, and $p=2^{m+1}$. Set
\begin{equation}\label{eq:witness}
 q=\max\left(0,\left\lceil\frac{M-r}{p}\right\rceil\right),
 \qquad j=r+pq.
\end{equation}
Then $j\geq M$ and
$j=2^m(2q+1)-1\in B_m\subseteq T_m$. Thus both displayed implications
apply to this particular integer. In fact $j$ is the least element of
$B_m$ above both cutoffs.

For example, with $m=3$, $N_0=100$, and $N_1=150$, formula
\eqref{eq:witness} gives $r=7$, $p=16$, $q=9$, and $j=151$.
The included program \texttt{obstruction\_witness.py} computes these
certificates using arbitrary-precision integers:
\begin{verbatim}
python3 obstruction_witness.py 3 100 150
\end{verbatim}
The program cannot infer cofinite bounds from an arbitrary black-box
infinite sequence. It verifies the final explicit implication once such
bounds are supplied. The proof of their existence under the amalgamation
hypothesis is mathematical, not computational.

\section{Separatedness and finite patching}
\label{sec:finite}

\subsection{Local equality always determines global equality}

\begin{theorem}[Separatedness, $\ZF$]\label{thm:separated}
For every set $A$, $\RA$ is separated for the dense topology. More
explicitly, if $D$ covers $F$ and $[f]_G=[g]_G$ for every $G\in D$,
then $[f]_F=[g]_F$.
\end{theorem}
\begin{proof}
Put $E=\eqset{f}{g}$. If $E\notin F$, the set $E^\co$ is
$F$-positive, and $H=\gen{F\cup\{E^\co\}}$ is proper. By density
there is $K\supseteq H$ in $D$. Local equality gives $E\in K$, while
$E^\co\in K$ follows from $H\subseteq K$. Their empty intersection
contradicts properness. Hence $E\in F$.
\end{proof}

This is the equality instance of the local-character phenomenon also
present in the source construction \cite[Lemma 5.8]{Massas2023}. The new
obstruction concerns existence of a global section, not failure of this
uniqueness property.

\subsection{Every finite compatible system has a representative}
The next theorem is slightly more general than needed: density is not
required for existence of a simultaneous representative.

\begin{theorem}[Finite pasting, $\ZF$]\label{thm:finite}
Let $A$ be nonempty, and let $F_1,\ldots,F_r$ be proper filters on an
index set $I$, where $1\leq r<\om$. Let $s_i\in A^I/{\sim_{F_i}}$.
Suppose the restrictions of $s_i$ and $s_j$ agree at every proper common
extension of $F_i$ and $F_j$. Then some $f:I\to A$ represents $s_i$
modulo $F_i$ for all $i$.
\end{theorem}
\begin{proof}
Choose representatives $f_i$ for the finitely many sections. For each
pair $i<j$ we claim there are $C_{ij}\in F_i$ and $C_{ji}\in F_j$
such that
\begin{equation}\label{eq:pairwitness}
 C_{ij}\cap C_{ji}\subseteq E(f_i,f_j).
\end{equation}
If $F_i,F_j$ are incompatible, choose their members with empty
intersection. If they are compatible, their generated common filter
$H_{ij}=\gen{F_i\cup F_j}$ is proper. Compatibility of the sections
implies $E(f_i,f_j)\in H_{ij}$. By the description of a generated
filter, it contains an intersection of one member of $F_i$ and one
member of $F_j$. This gives \eqref{eq:pairwitness}. All the selections
just made are finite and are available in $\ZF$.

For each $i$, intersect the finitely many $C_{ij}$ on its side and call
the result $C_i\in F_i$. For $r=1$, take $C_1=I$. Then
\[
 C_i\cap C_j\subseteq E(f_i,f_j)\qquad(i\neq j).
\]
Fix $a\in A$ and define $f(t)$ as follows. If $t$ belongs to at least
one $C_i$, use $f_i(t)$ for the least such $i$; otherwise use $a$.
Whenever $t\in C_i$, the chosen value agrees with $f_i(t)$ by the
intersection property. Thus $C_i\subseteq E(f,f_i)$, so
$[f]_{F_i}=s_i$ for every $i$.
\end{proof}

\begin{corollary}[Finite covers and sharpness]\label{cor:sharp}
Every matching family for $\RA$ on a finitely generated covering sieve
has a unique amalgam. For $|A|\geq2$, there is a failure of amalgamation
on a countably generated covering sieve. Hence a witness cannot have
finitely many generators, while countably many suffice.
\end{corollary}
\begin{proof}
Apply Theorem~\ref{thm:finite} to sections at the finite generating
filters. Their restrictions determine the given matching family on the
entire sieve. The resulting representative yields a section at the base
filter, and Theorem~\ref{thm:separated} gives uniqueness. If $A$ is
empty there is no matching family on a nonempty cover. The second
assertion is the construction in Section~\ref{sec:obstruction}.
\end{proof}

The word ``countable'' in this corollary refers to the number of
\emph{sieve generators}, not the number of filters in the full covering
sieve. In particular, the obstruction does not contradict finite pasting:
any finite part of its prescribed data can be represented, but the whole
cover cannot.

\section{The obstruction above an arbitrary base filter}
\label{sec:antichain}

The partition proof extends whenever a base filter admits infinitely many
disjoint positive sets. This is the local form that will support the later
classification.

\begin{theorem}[Positive-row obstruction, $\ZF$]\label{thm:positive-rows}
Let $F\in\cF$. Suppose $(B_n)_{n<\om}$ is a sequence of pairwise
disjoint $F$-positive subsets of $\om$. They need not cover $\om$.
For every set $A$ with two distinct elements, $\RA$ fails the sheaf
condition at $F$ on the site of all free filters extending $F$.
\end{theorem}
\begin{proof}
Define
\[
 T_m=\om\setminus\bigcup_{n<m}B_n,\qquad
 H_n=\gen{F\cup\{B_n\}},\qquad
 H_\infty=\gen{F\cup\{T_m:m<\om\}}.
\]
Each $H_n$ is proper by positivity. For the tail filter, any finite
intersection of its added generators is some $T_m$. If $C\in F$,
then $C\cap T_m$ contains $C\cap B_m$, which is infinite. Thus
$H_\infty$ is proper as well. Notice its precise membership test:
\begin{equation}\label{eq:general-tail-membership}
 X\in H_\infty
 \quad\Longleftrightarrow\quad
 \text{some }C\in F\text{ and }m<\om\text{ satisfy }C\cap T_m\subseteq X.
\end{equation}
The row filters are pairwise incompatible, and $H_n$ is incompatible
with $H_\infty$ because $B_n\cap T_{n+1}=\varnothing$.

The sieve
\[
 D_F=\uparrow H_\infty\ \cup\ \bigcup_n\uparrow H_n
\]
is dense above $F$. Given $G\supseteq F$, a $G$-positive row may be
adjoined properly. If none is positive, all row complements lie in $G$;
therefore all $T_m$ lie in $G$, and $G\supseteq H_\infty$.

Assign $[\const{a_0}]$ on the row components and $[\const{a_1}]$ on
the tail component. This is again a matching family. If $[f]_F$ glues
it, let $E_i=\{j:f(j)=a_i\}$. Membership $E_1\in H_\infty$ gives
$C\in F$ and $m$ with $C\cap T_m\subseteq E_1$.
Membership $E_0\in H_m$ gives $C'\in F$ with
$C'\cap B_m\subseteq E_0$. Since $B_m$ is $F$-positive,
\[
 (C\cap C')\cap B_m
\]
is infinite and hence nonempty. It is contained in both $E_0$ and
$E_1$, a contradiction.
\end{proof}

\begin{remark}
Using $T_m=\bigcup_{n\geq m}B_n$ in this more general proof would be
unjustified when the rows do not cover $\om$: avoiding every row would
not force membership of that union. The complement-of-a-finite-union
definition is the correct one and includes any leftover points.
\end{remark}

\section{Every countably generated free base already fails}
\label{sec:countablebase}

The main theorem uses a particularly simple base. In fact, no countably
generated free base can repair the gluing problem.

\begin{lemma}[Disjoint positive rows from a countable base, $\ZF$]
\label{lem:countable-rows}
Every countably generated proper free filter on $\om$ has a sequence
of pairwise disjoint positive sets.
\end{lemma}
\begin{proof}
Fix a countable base and replace it by successive finite intersections
to obtain a decreasing base
\[
 C_0\supseteq C_1\supseteq C_2\supseteq\cdots.
\]
Every $C_s$ is infinite. At stage $s$, for $n=0,\ldots,s$, let
$b_{n,s}$ be the least element of $C_s$ not used previously, including
earlier choices within that stage. Only finitely many numbers have
been used at this point, so the least available element exists.
This is a recursion with a uniquely determined value at every step,
not an invocation of dependent choice.

Let $B_n=\{b_{n,s}:s\geq n\}$. No chosen integer is ever reused,
so these sets are pairwise disjoint. If $X\in F$, some $C_m$ is
contained in $X$. For every $s\geq\max(m,n)$,
\[
 b_{n,s}\in C_s\subseteq C_m\subseteq X.
\]
These are infinitely many distinct elements of $B_n\cap X$.
Lemma~\ref{lem:positive} proves that $B_n$ is $F$-positive.
\end{proof}

\begin{corollary}[$\ZF$]\label{cor:countablebase}
If $F$ is a countably generated proper free filter and $|A|\geq2$,
then $\RA$ fails sheafness at $F$ on a countably generated covering
sieve. Its generating filters can also all be taken countably generated.
\end{corollary}
\begin{proof}
Apply Theorem~\ref{thm:positive-rows} to the preceding lemma. Adding
one set, or a countable decreasing family of sets, to a countable filter
base leaves a countable base, provided the generated filter is proper.
Properness was proved in Theorem~\ref{thm:positive-rows}.
\end{proof}

\subsection{Restricting the site does not remove this obstruction}
Let $\cF_{\mathrm{cg}}$ be the category consisting only of countably
generated proper free filters, with the same inclusion arrows and the
relative dense topology. For an arbitrary $G$ in this smaller category,
adjoining a positive row still gives a countably generated proper
extension. The other branch of the density argument uses $G$ itself.
Thus the sieve of extensions of the displayed generators \emph{within}
$\cF_{\mathrm{cg}}$ is still covering, and the same family still fails
to glue. This applies above every object of that category.

Consequently, the main phenomenon is not caused by admitting complicated
uncountably generated filters or by the presence of ultrafilters.
The elementary countable-base part of the site already exhibits it.
This observation does not assert that every free filter is countably
generated; no such assertion is used.

\section{A local Boolean-algebra classification}
\label{sec:boolean}

\subsection{The quotient attached to a base}
Fix $F\in\cF$. Define its dual ideal and quotient Boolean algebra by
\begin{equation}\label{eq:BF}
 \I_F=\{S\subseteq\om:S^\co\in F\},\qquad
 \B_F=\Pow(\om)/\I_F.
\end{equation}
Let $q_F:\Pow(\om)\to\B_F$ be the quotient map. The ideal is proper
and contains every finite set. We have
\begin{equation}\label{eq:quotient-tests}
 q_F(S)=0\Longleftrightarrow S^\co\in F,
 \qquad
 q_F(S)=1\Longleftrightarrow S\in F.
\end{equation}
Thus $F$-positive sets represent exactly the nonzero Boolean elements.

\begin{lemma}[Filters above $F$]\label{lem:filter-correspondence}
Taking image and inverse image under $q_F$ gives a bijection between
proper filters on $\om$ extending $F$ and proper filters of $\B_F$.
\end{lemma}
\begin{proof}
If $G\supseteq F$ and $S\triangle T\in\I_F$, then
$(S\triangle T)^\co\in F\subseteq G$. Whenever $S\in G$,
intersecting $S$ with this complement gives a member of $G$ contained
in $T$, so $T\in G$. Hence membership in $G$ depends only on the
quotient class. Its image is therefore a proper Boolean filter.
Conversely, the inverse image of a proper Boolean filter is a proper
filter on $\om$ and contains $q_F^{-1}(1)=F$. These operations
are inverse by the same membership test.
\end{proof}

\subsection{The finite case needs no choice}

\begin{proposition}[Finite quotient implies sheafness, $\ZF$]
\label{prop:finiteB}
If $\B_F$ is finite, then for every set $A$, the restriction of
$\RA$ to the category of filters extending $F$ is a sheaf.
\end{proposition}
\begin{proof}
A finite Boolean algebra has only finitely many proper filters.
Lemma~\ref{lem:filter-correspondence} shows that the category above
$F$ has finitely many objects. Thus any covering sieve at any of its
objects is generated by finitely many filters, namely its own members.
Corollary~\ref{cor:sharp} proves the sheaf assertion. Empty value sets
are covered by the same vacuity argument as in Theorem~\ref{thm:main}.
\end{proof}

There is a useful concrete description of this case. If $\B_F$ has
atoms $b_1,\ldots,b_k$, let
\[
 U_i=\{X\subseteq\om:b_i\leq q_F(X)\}.
\]
Each $U_i$ is an ultrafilter extending $F$: an atom lies below exactly
one of a Boolean element and its complement. These are all the
ultrafilter extensions of $F$, and
\begin{equation}\label{eq:finite-intersection}
 F=\bigcap_{i=1}^kU_i.
\end{equation}
No ultrafilter lemma is needed to construct these finitely many filters.

Every proper filter of a finite Boolean algebra is generated by its least
nonzero member. Therefore every $G\supseteq F$ has the form
$G=\bigcap_{i\in J}U_i$ for a nonempty subset $J\subseteq\{1,\ldots,k\}$.
For nonempty $A$ the natural map is a bijection:
\begin{equation}\label{eq:finiteproduct}
 A^\om/G\ \longrightarrow\ \prod_{i\in J} A^\om/U_i,
 \qquad [f]_G\longmapsto([f]_{U_i})_{i\in J}.
\end{equation}
Injectivity follows from the expression for $G$. For surjectivity, take
representatives for the finitely many given classes and paste them on
disjoint representatives of the Boolean atoms. Such disjoint
representatives exist by finite disjointification; the negligible leftover
set can be assigned to the first piece. The resulting sequence has the
prescribed class at each relevant ultrafilter. This is another instance
of finite pasting. For $A=\varnothing$ both sides are empty since
$J$ is nonempty.

\subsection{The infinite case under dependent choice}

\begin{lemma}[$\ZF+\DC$]\label{lem:BA-antichain}
Every infinite Boolean algebra has a countably infinite sequence of
pairwise disjoint nonzero elements.
\end{lemma}
\begin{proof}
For an element $r$, write $\B\mathbin{\upharpoonright}r$ for the
interval of elements below $r$. Start with $r_0=1$, whose interval
is infinite. An element $r$ with infinite interval is not an atom, so
it can be split into two disjoint nonzero elements $c,d$ with
$c\vee d=r$. At least one of the intervals below $c,d$ is infinite:
otherwise the map $x\mapsto(x\wedge c,x\wedge d)$ would embed the
interval below $r$ into a finite product of finite sets.

Use dependent choice to repeat this operation, retaining an infinite
interval as $r_{n+1}$ and calling the other nonzero piece $b_n$.
Then $b_n\wedge r_{n+1}=0$ and all later $b_j$ lie below
$r_{n+1}$. Hence the $b_n$ are nonzero and pairwise disjoint.
\end{proof}

\begin{lemma}[$\ZF+\DC$]\label{lem:lift-antichain}
If $\B_F$ is infinite, there is a sequence of literally disjoint
$F$-positive subsets of $\om$.
\end{lemma}
\begin{proof}
Take the sequence $(b_n)$ supplied by Lemma~\ref{lem:BA-antichain}.
Dependent choice implies countable choice, so select representatives
$C_n\subseteq\om$ with $q_F(C_n)=b_n$. Define
\[
 B_n=C_n\setminus\bigcup_{j<n}C_j.
\]
These sets are literally disjoint. Because $b_n\wedge b_j=0$ for
$j<n$, their quotient classes satisfy
\[
 q_F(B_n)=b_n\wedge\bigwedge_{j<n}\neg b_j=b_n\neq0.
\]
The last inequality is exactly $F$-positivity by
\eqref{eq:quotient-tests}.
\end{proof}

\begin{theorem}[Local classification, $\ZF+\DC$]\label{thm:local}
Let $F\in\cF$ and let $A$ have at least two elements. The following
conditions are equivalent:
\begin{enumerate}[label=\textup{(\roman*)}]
 \item $\RA$ is a sheaf on the full category of filters extending $F$,
       equipped with its dense topology;
 \item every matching family on every covering sieve at $F$ has an amalgam;
 \item the Boolean algebra $\B_F$ is finite.
\end{enumerate}
The implication from \textup{(iii)} to \textup{(i)} is provable in $\ZF$.
\end{theorem}
\begin{proof}
The first condition implies the second by definition, and
Proposition~\ref{prop:finiteB} gives the third-to-first implication.
If $\B_F$ is infinite, Lemma~\ref{lem:lift-antichain} and
Theorem~\ref{thm:positive-rows} produce a matching family at $F$
without an amalgam. This is the contrapositive of the remaining
implication.
\end{proof}

For arbitrary $A$, the classification may be written
\[
 \RA\text{ is a sheaf above }F
 \quad\Longleftrightarrow\quad
 |A|\leq1\ \text{ or }\ \B_F\text{ is finite},
 \qquad(\ZF+\DC).
\]
For $F=\Fr$, the quotient is $\Pow(\om)/\mathrm{Fin}$ and the
explicit rows of Section~\ref{sec:arithmetic} already give an infinite
antichain, so dependent choice is unnecessary there.

The local theorem also identifies genuine positive cases. If a free
ultrafilter $U$ is given, the category above $U$ has a single object;
$\B_U$ is the two-element Boolean algebra, and every $\RA$ is a
sheaf on that one-object site. This does not conflict with failure on the
larger category of all free filters.

\section{What sheafification adds in classical set theory}
\label{sec:sheafification}

\subsection{The product of all ultrafilter values}
This section assumes $\ZFC$. We make the associated sheaf completely
explicit, rather than relying on an abstract sheafification construction.
For a free filter $F$, let
\[
 \Ult(F)=\{U:U\text{ is an ultrafilter on }\om\text{ and }F\subseteq U\}.
\]
Define
\begin{equation}\label{eq:sheafification}
 \SA(F)=\prod_{U\in\Ult(F)} A^\om/U.
\end{equation}
When $G\supseteq F$, its ultrafilter extensions form a subset of
$\Ult(F)$, and restriction in $\SA$ is the corresponding projection.
There is a natural transformation
\begin{equation}\label{eq:eta}
 \eta_F:\RA(F)\longrightarrow\SA(F),\qquad
 [f]_F\longmapsto([f]_U)_{U\in\Ult(F)}.
\end{equation}

\begin{lemma}\label{lem:ultra-cover}
The ultrafilters extending $F$ form a covering sieve at $F$.
Every covering sieve at $F$ contains every member of $\Ult(F)$.
Each map $\eta_F$ is injective.
\end{lemma}
\begin{proof}
Every proper filter extends to an ultrafilter by the ultrafilter lemma,
which is available in $\ZFC$. Therefore these maximal extensions are
dense. A proper extension of an ultrafilter equals itself, so this dense
family is extension-closed and hence a sieve.

If $D$ is any covering sieve and $U\in\Ult(F)$, density applied to
$U$ gives a member of $D$ extending $U$. It must be $U$ itself.

For injectivity, suppose $[f]_F\neq[g]_F$. Then
$E(f,g)\notin F$, so $\gen{F\cup\{E(f,g)^\co\}}$ is proper.
Extend it to an ultrafilter $U$. Its containing $E(f,g)^\co$ prevents
$E(f,g)\in U$, and the $U$-coordinates of the two images differ.
\end{proof}

\begin{theorem}[Explicit associated sheaf, $\ZFC$]\label{thm:sheafification}
The presheaf $\SA$ is a sheaf, and $\eta:\RA\to\SA$ has the
universal property of sheafification.
\end{theorem}
\begin{proof}
First let $(s_G)_{G\in D}$ be a matching family for $\SA$ on a cover
of $F$. For each $U\in\Ult(F)$, Lemma~\ref{lem:ultra-cover} gives
$U\in D$. Since $U$ has only itself as an ultrafilter extension,
$\SA(U)=A^\om/U$. Let the $U$-coordinate of $s\in\SA(F)$ be
$s_U$. For $G\in D$ and $U\supseteq G$, compatibility says that
the $U$-coordinate of $s_G$ is $s_U$. Hence $s$ restricts to $s_G$.
The same coordinates show uniqueness. This proves sheafness.

Now let $\mathscr T$ be any sheaf and
$\varphi:\RA\to\mathscr T$ any morphism of presheaves. Given
$s\in\SA(F)$, assign the section
\[
 \varphi_U(s(U))\in\mathscr T(U)
 \quad\text{to each }U\in\Ult(F).
\]
Distinct ultrafilters are incompatible, so these sections form a
matching family on the ultrafilter cover. Glue them uniquely in
$\mathscr T(F)$ and call the result $\psi_F(s)$. These definitions
commute with restriction: both possible restrictions at a stronger
filter have the same values at all its ultrafilter extensions and are
therefore equal by sheaf uniqueness. Thus $\psi$ is a morphism.

For $s=\eta_F([f]_F)$, naturality of $\varphi$ shows that the
sections being glued are the restrictions of $\varphi_F([f]_F)$.
Hence $\psi\eta=\varphi$. Conversely, any morphism with this
property must take $s$ to a section with these prescribed ultrafilter
restrictions; sheaf uniqueness forces that section to be $\psi_F(s)$.
This proves the universal property.
\end{proof}

This proof also explains why sheafification does not collapse unequal
old sections: every $\eta_F$ is injective. It supplies additional
sections by allowing independent values at maximal filters, not by
identifying distinct equivalence classes already present.

\subsection{The two-element case: clopen sets versus arbitrary sets}
Let $A=2=\{0,1\}$. Every sequence $f:\om\to2$ is a characteristic
function. At an ultrafilter $U$, exactly one of its two fibres lies in
$U$, so $2^\om/U$ is canonically $2$. Accordingly,
\begin{equation}\label{eq:binary-product}
 \mathscr S_2(F)=2^{\Ult(F)}\cong\Pow(\Ult(F)),
 \qquad \mathscr R_2(F)\cong\B_F.
\end{equation}
The second identification sends a characteristic function to the
quotient class of its support; equality of functions modulo $F$ means
that the symmetric difference of their supports belongs to $\I_F$.

For $X\subseteq\om$, set
\[
 \widehat X=\{U\in\Ult(F):X\in U\}.
\]
These sets give the usual Stone topology on $\Ult(F)$: each is both
open and closed, finite intersections and complements have the expected
Boolean formulas, and they form a base. For completeness, compactness
follows directly from the ultrafilter lemma. A family of these basic
sets with no finite subcover has complements whose representatives,
together with $F$, have the finite-intersection property; an ultrafilter
extending that family is missed by the alleged cover. A general open
cover can be refined by the basic sets. Thus this space is compact.
Every clopen subset, being compact and open, is a finite union of basic
sets and is therefore itself $\widehat X$ for some $X$.

Under \eqref{eq:binary-product}, $\eta_F$ is therefore the inclusion
\begin{equation}\label{eq:stone}
 \B_F\cong\operatorname{Clop}(\Ult(F))
 \ \hookrightarrow\ \Pow(\Ult(F)).
\end{equation}
This is a description of the sheafification map, not a claim that the
power set on the right is the Boolean completion in the order-dense or
MacNeille sense. Those are different requirements.

For any finite nonempty $A$, the same description reads: the old
sections induce the continuous maps from the Stone space to discrete
$A$, while the associated sheaf permits all maps. Each fibre of the
map induced by a sequence is a basic clopen set. Conversely, a
continuous map gives a finite clopen partition; representing its pieces
in $\Pow(\om)$ and removing the finitely many negligible overlaps
produces a sequence with precisely those ultrafilter values.

\subsection{The missing section in the explicit example}
Return to $F=\Fr$ and the rows and tails of
Section~\ref{sec:obstruction}. Define
\begin{equation}\label{eq:K}
 K=\{U\in\Ult(\Fr):F_\infty\subseteq U\}
  =\bigcap_m\widehat{T_m}.
\end{equation}
Every free ultrafilter contains at most one row. If it contains no row,
it contains the complement of each row and therefore all $T_m$.
Hence every free ultrafilter either extends a unique $F_n$ or belongs
to $K$. The two-valued matching family becomes the function
\[
 e(U)=
 \begin{cases}
  0,&U\supseteq F_n\text{ for some }n,\\
  1,&U\in K.
 \end{cases}
\]
It is the characteristic function of $K$ in $\mathscr S_2(\Fr)$.

\begin{proposition}\label{prop:K}
The set $K$ is closed but not clopen. Equivalently, the section $e$
is not in the image of $\eta_{\Fr}$ for $A=2$.
\end{proposition}
\begin{proof}
Equation~\eqref{eq:K} makes $K$ closed. Suppose $K=\widehat E$ for
some $E\subseteq\om$. Then $E$ lies in every ultrafilter extending
$F_\infty$, which forces $E\in F_\infty$: otherwise one could
extend $\gen{F_\infty\cup\{E^\co\}}$ to an ultrafilter. Thus
$T_m\setminus E$ is finite for some $m$.

Take an ultrafilter $U\supseteq F_m$, which exists in $\ZFC$.
Since $B_m\subseteq T_m$ and $B_m\setminus E$ is finite, we have
$E\in F_m\subseteq U$. Hence $U\in\widehat E=K$. But no
proper filter extends both $F_m$ and $F_\infty$, a contradiction.
Thus $K$ is not clopen.
\end{proof}

Closed but nonclopen does not imply that $K$ has empty interior. Indeed,
one can choose the explicit selector $S=\{2^n-1:n<\om\}$, which
meets each row once and is almost contained in every $T_m$.
The nonempty basic open set $\widehat S$ is contained in $K$.
This observation helps avoid a misleading picture of the tail component
as necessarily nowhere dense.

\section{Consequences for the hyperreal setting}
\label{sec:hyperreals}

When $A=\mathbb R$, pointwise ring operations make the reduced powers
and their restriction maps into rings and ring homomorphisms. At an
ultrafilter $U$ the reduced power is the usual ultrapower; over the
whole filter site it is a varying family, not one fixed field. The
underlying set obstruction uses only $0,1\in\mathbb R$, so no
ordered-field property, transfer theorem, saturation theorem, or
cardinal estimate for $\mathbb R$ is required.

In the product sheaf of Section~\ref{sec:sheafification}, the same
coordinatewise $0$--$1$ element $e$ satisfies $e^2=e$. It does not come
from any sequence of reals. To see this directly, if $\eta([f])=e$,
then at every row filter the injectivity in
Lemma~\ref{lem:ultra-cover} forces $[f]_{F_n}=[0]_{F_n}$, and at
the tail filter it forces $[f]_{F_\infty}=[1]_{F_\infty}$.
Proposition~\ref{prop:contradiction} rules this out.

This is a \emph{new global idempotent section relative to the original
presheaf}, not a claim that its old rings had no nontrivial idempotents.
For example, characteristic sequences already give idempotents in
$\mathbb R^\om/\Fr$. Nor is the new global section a nontrivial
idempotent inside one hyperreal field: at each ultrafilter it is just
$0$ or $1$. The product of many fields is not itself a field.

Three statements should therefore be kept separate. Local equality of
old sections detects equality at the base filter; this is separatedness.
Compatible local sections need not have a single representative; this
is the obstruction. Passing to the associated sheaf supplies such
sections but changes the available global elements. None of these
statements contradicts the possibility-semantic account in the source,
which already distinguishes its presheaf construction from a sheaf
model \cite[Section 5.2]{Massas2023}.

The result also depends on the category and its topology. It is not an
assertion that every sheaf-based construction of nonstandard analysis
fails, or that every reduced-power presheaf on a different filter site
has the same behavior. The relative countably generated site considered
in Section~\ref{sec:countablebase} was analyzed separately; no
unexamined change of topology is implicit in the main theorem.

\section{Reproducibility and proof audit}
\label{sec:audit}

\subsection{Exact supplementary checks}
The mathematical proof is infinite and does not rest on numerical
experiments. The archive nevertheless supplies executable checks of the
arithmetic identities and of a finite analogue of the pasting theorem.
All code uses the Python standard library and exact integers.
Running
\begin{verbatim}
python3 verify.py --output verification_results.json
\end{verbatim}
reproduces the following results:
\begin{center}
\begin{tabular}{lr}
\toprule
Check & Number checked\\
\midrule
Row--tail membership identities & 150,000\\
Exact cutoff-to-witness certificates & 64,925\\
Invalid inputs correctly rejected & 6\\
Finite generator families & 717\\
Finite local configurations & 86,669\\
Compatible finite configurations & 9,209\\
Unique gluings for finite covers & 7,421\\
Compatible noncover pastings & 1,788\\
\bottomrule
\end{tabular}
\end{center}
All checks passed. Of the cutoff certificates, $54,925$ exhaust
$m=0,\ldots,12$ and $N_0,N_1=0,\ldots,64$; another $10,000$ use a
fixed pseudorandom seed, $m\leq200$, and $512$-bit nonnegative
cutoffs. The program also checks that the returned witness is the
least possible one.

For the finite-pasting tests, a nonempty subset $K$ of a finite index
set $I$ determines the principal filter
$F_K=\{X\subseteq I:K\subseteq X\}$. A reduced-power class is
exactly a function on $K$. The program exhausts families of at most
three distinct nonempty kernels, all local functions, and all global
functions for $|I|\leq4$, $|A|=2$, and $|I|\leq3$, $|A|=3$.
Compatibility is checked on pairwise intersections. The least-index
pasting construction is then compared with an independent enumeration
of global solutions. A family covers the weakest proper filter exactly
when the union of its kernels is $I$; in that case uniqueness is checked.

A finite index set has no proper filter containing all of its cofinite
subsets, since every subset, including the empty set, is cofinite there.
These finite tests are therefore \emph{not} finite objects of the
Fr\'echet site. They test the more general finite-pasting theorem of
Section~\ref{sec:finite}, which was deliberately stated for arbitrary
proper filters on an index set. The tests cannot establish sheaf failure
on an infinite site, and no proof-assistant verification is claimed.

\subsection{Foundational dependencies}
\begin{center}
\small
\begin{tabular}{p{0.64\textwidth}p{0.23\textwidth}}
\toprule
Result & Assumptions used\\
\midrule
Explicit counterexample and global classification & $\ZF$\\
Separatedness and finite-pasting theorem & $\ZF$\\
Positive-row obstruction at a fixed base & $\ZF$\\
Failure at every countably generated free base & $\ZF$\\
Finite Boolean quotient implies local sheafness & $\ZF$\\
Infinite Boolean quotient gives positive rows & $\ZF+\DC$\\
Product-ultrapower associated sheaf; Stone description & $\ZFC$\\
\bottomrule
\end{tabular}
\end{center}
The stronger assumptions in the final rows are not smuggled back into
the first rows. In particular, the main construction does not choose any
ultrafilter. The $\ZFC$ argument is an interpretation and extension of a
result already established in $\ZF$.

\subsection{The main argument as a checklist}
The decisive checks are all explicit. The row and tail families are
proper filters; their incompatibility follows from displayed disjoint
members. The cover includes \emph{all} extensions of each generator.
For every test filter, a positive row can be adjoined, or all finite row
tails already belong to it. Each local section is the class of a
specified constant sequence. Incompatibility makes the assignment
unambiguous and persistent. Finally, membership at the tail filter
produces one tail index, and membership at that row filter produces
the contradiction on an infinite set.

No step requires that a matching family be globally represented before
trying to glue it. No countable intersection is asserted to belong to
a filter. The empty-domain convention is handled explicitly. These are
the main points at which an apparently similar argument could fail;
they are also the points against which the proof was checked.

\section{Conclusion and the remaining status questions}
\label{sec:conclusion}

The selected published question has the following answer within the
argument of this manuscript: uncountability is not necessary. Two values
suffice, and a value set has a sheaf reduced-power presheaf on the full
Fr\'echet site exactly when it has at most one element. The proof uses an
explicit row partition and a tail filter and is available in $\ZF$.

The extensions identify the mechanism rather than merely producing an
isolated counterexample. Finite systems always paste. Countably generated
free bases always admit the obstruction. With dependent choice, the
finite or infinite nature of the quotient Boolean algebra completely
determines local sheafness for a nontrivial value set. In classical set
theory, the associated sheaf replaces the requirement of a single
sequence by a compatible product of ultrapower values, and the missing
section is already binary.

The mathematical arguments are presented as complete proofs, but this
manuscript has not undergone independent expert review or formal
verification. The literature check establishes a documented open
question in 2023 and its repetition in 2024; it does not establish that
no later or unpublished proof exists. Those limitations concern
certification and priority, not an omitted step attributed to the
problem's age or difficulty.

A precise further foundational question suggested by the proofs is
whether the infinite-Boolean-quotient direction of
Theorem~\ref{thm:local} can be established in $\ZF$ without additional
hypotheses on the base filter. This manuscript does not assert that the
use of dependent choice there is optimal, nor that this suggested
question is new or open in the literature. It records exactly where
choice enters the present extension.

\appendix
\section{A compact standalone proof}
\label{app:shortproof}
Let $\om=\bigsqcup_nB_n$ be a partition into infinite sets and
$T_m=\bigcup_{n\geq m}B_n$. Let $F_n$ consist of sets containing
all but finitely many elements of $B_n$, and let $F_\infty$ consist
of sets containing all but finitely many elements of some $T_m$.
All are proper free filters. Distinct $F_n$ are incompatible, and
$F_n$ is incompatible with $F_\infty$ because
$B_n\cap T_{n+1}=\varnothing$.

The sieve of all extensions of these filters covers $\Fr$. Indeed,
if a free filter $G$ has a positive row, adjoining it gives an
extension in the sieve. Otherwise $G$ contains every row complement,
hence every $T_m$, so $G$ itself extends $F_\infty$.

For distinct $a_0,a_1\in A$, assign the constant $a_0$ class on row
components and the constant $a_1$ class on the tail component.
Incompatibility makes this a matching family. A gluing sequence $f$
would equal $a_1$ cofinitely on some $T_m$ and equal $a_0$ cofinitely
on $B_m$. Since $B_m\subseteq T_m$ is infinite, this is impossible.
For $|A|\leq1$, the sheaf condition is immediate. This proves the
classification without any choice axiom.

\section{Source map and archive contents}
\label{app:sources}
The exact source locations are as follows. In \cite{Massas2023}, the
filter category, covers, and reduced-power presheaf are Definitions
5.4, 5.5, and 5.7 on printed page~515. The sufficient cardinality
condition is Fact~5.10 on pages~516--517. The motivating question is
the paragraph immediately after its proof on page~517. The 2024
version is \cite[Section 7.4.2, Fact 7.4.10 and following paragraph]{Massas2024},
with the question on printed page~249. The PDF page indices are
respectively $27$ and $259$ when counted from zero. These passages
were checked against page images as well as extracted text.

The archive contains \texttt{article.tex}, \texttt{article.pdf}, the
bibliographic database \texttt{references.bib}, the two Python
programs, their actual \texttt{verification\_results.json}, a
\texttt{README.md}, a separate proof audit and literature-search record,
and a build script. Source publications and font files are not bundled.
The PDF bibliography is self-contained; the BibTeX database is supplied
for reuse and is not required by the build.

\begin{thebibliography}{9}
\addcontentsline{toc}{section}{References}
\bibitem{Massas2023}
Guillaume Massas.
\newblock A Semi-Constructive Approach to the Hyperreal Line.
\newblock \emph{Australasian Journal of Logic} \textbf{20}(3) (2023),
490--536, article~6.
\newblock DOI: \href{https://doi.org/10.26686/ajl.v20i3.8254}{10.26686/ajl.v20i3.8254}.
\newblock \url{https://ojs.victoria.ac.nz/ajl/article/view/8254}.
\newblock Open question: p.~517, after Fact~5.10.

\bibitem{Massas2024}
Guillaume Massas.
\newblock \emph{Duality and Infinity}.
\newblock Ph.D. dissertation, University of California, Berkeley, 2024.
\newblock \url{https://escholarship.org/uc/item/8pw910s3}.
\newblock Chapter~7, Section~7.4.2; open question: p.~249, after
Fact~7.4.10.

\bibitem{StacksTopology}
The Stacks Project Authors.
\newblock The topology defined by a site.
\newblock \emph{The Stacks Project}, Section~7.48, tag~00ZB.
\newblock \url{https://stacks.math.columbia.edu/tag/00ZB}.
\newblock Consulted 20 September 2026.
\end{thebibliography}

\end{document}
